\section{System architecture}\label{sec:arch}
Figure~\ref{fig:arch} shows the inference path. A sweep of each breast is a sequence of frames. Every frame is preprocessed (overlay removal, field-of-view crop, intensity normalisation, letterboxing to $224\times224$) and scored by the quality gate~Q. The frozen encoder produces a CLS token and a $16\times16$ grid of patch tokens. The detector~D maps patch tokens to a lesion probability map, from which it derives a frame lesion probability, candidate boxes and a heatmap. The malignancy head~M scores the top box. The aggregator~A requires a lesion to persist over consecutive adequate frames, takes the maximum over each breast, and combines the two breasts and their coverage into the woman-level decision. The operating threshold comes from the safety layer, which calibrates scores and certifies the threshold at the deployment site. Every request is logged with the model version.

\begin{figure*}[t]
\centering
\begin{tikzpicture}[
  node distance=5mm and 4mm,
  box/.style={draw=line, rounded corners=2pt, fill=white, align=center, font=\footnotesize, minimum height=12mm, inner sep=2pt, text width=20mm},
  key/.style={box, draw=acc, line width=0.8pt},
  decision/.style={box, draw=acc2, line width=0.8pt, text width=21mm},
  arr/.style={-{Stealth[length=2mm]}, draw=ink2, line width=0.6pt},
  lab/.style={font=\scriptsize, text=ink2}]
\node[box] (sweep) {Handheld sweep\\left + right\\breast frames};
\node[box, right=of sweep] (prep) {Preprocessing\\overlay inpaint,\\FOV crop, norm.,\\letterbox 224};
\node[key, right=of prep] (q) {Q: quality gate\\frame adequacy,\\breast coverage};
\node[box, right=of q] (enc) {Encoder\\DINOv2 ViT-S/14\\(frozen), CLS +\\$16{\times}16$ tokens};
\node[key, right=of enc] (det) {D: detector\\patch head $\rightarrow$\\frame prob.,\\boxes, heatmap};
\node[key, right=of det] (agg) {A: aggregation\\persistence $\rightarrow$\\breast $\rightarrow$\\woman};
\node[decision, right=of agg] (dec) {\textsc{normal}\\\textsc{recall}\\\textsc{unreadable}};
\node[key, below=of det] (mal) {M: malignancy\\head, CLS +\\ROI tokens\\(+ age)};
\node[key, below=of agg] (safe) {Safety layer\\calibration,\\certified\\threshold};
\node[box, below=of dec] (log) {Service: API,\\web UI, model\\registry,\\JSONL log};
\draw[arr] (sweep) -- (prep);
\draw[arr] (prep) -- (q);
\draw[arr] (q) -- (enc);
\draw[arr] (enc) -- (det);
\draw[arr] (det) -- (agg);
\draw[arr] (agg) -- (dec);
\draw[arr] (det) -- (mal);
\draw[arr] (mal.north east) -- (agg.south west);
\draw[arr] (safe) -- node[lab, right] {threshold} (agg);
\draw[arr] (q.north) -- ++(0,6mm) -| node[lab, pos=0.25, above] {adequacy flags} (agg.north);
\draw[arr] (dec) -- (log);
\end{tikzpicture}
\caption{SweepScreen inference path. Learned components (blue outline) are fitted on public data with the shared training code used for every evaluation fold; the output (orange) is one decision per woman. M passes a suspicion score for each detected lesion to A, which reports it with every recall.}
\label{fig:arch}
\end{figure*}

All learned components are fitted by one function (\code{sweepscreen.train.fit}). Every evaluation fold and the deployed model use the same code path, so the reported numbers describe the code that is deployed.
